% 演習問題集——第2章
% 章番号はこのファイル名によって決まる。

\begin{exo}[低温の氷から沸騰水まで：オーダー評価]{chauffage-glace-eau-ebullition}
\seoready{true}
\keywords{熱量測定, 相転移, 融解, 気化, 潜熱, オーダー}

\textbf{数値データ（大気圧下）：}
\[
\begin{aligned}
c_{\text{ice}} &= 2{,}1\times10^3\,\text{J}\!\cdot\!\text{kg}^{-1}\!\cdot\!\text{K}^{-1},
& L_f &= 3{,}34\times10^5\,\text{J}\!\cdot\!\text{kg}^{-1},\\
c_{\text{water}} &= 4{,}18\times10^3\,\text{J}\!\cdot\!\text{kg}^{-1}\!\cdot\!\text{K}^{-1},
& L_v &= 2{,}26\times10^6\,\text{J}\!\cdot\!\text{kg}^{-1}.
\end{aligned}
\]
大気圧下で，初期温度 $-100\,{}^\circ\text{C}$ の氷1キログラムを段階的に加熱する。

\begin{enumerate}
  \item 氷を $-100\,{}^\circ\text{C}$ から $0\,{}^\circ\text{C}$ まで加熱するのに必要なエネルギーを求めよ。
  \item $0\,{}^\circ\text{C}$ で氷を完全に融解させるのに必要なエネルギーを求めよ。
  \item 液体の水を $0\,{}^\circ\text{C}$ から $100\,{}^\circ\text{C}$ まで加熱するのに必要なエネルギーを求めよ。
  \item $100\,{}^\circ\text{C}$ で水を完全に気化させるために追加で必要なエネルギーを求めよ。
  \item どの段階が最も多くのエネルギーを必要とするか。1キログラムの水を $0\,{}^\circ\text{C}$ から $100\,{}^\circ\text{C}$ まで加熱するエネルギーと，質量 $m$ が10メートル落下するときの重力位置エネルギーを比較せよ。同じエネルギーを放出するには，どれだけの質量を落下させればよいか。
  \item 逆過程を考える。質量 $m_v$ の水蒸気がフロントガラス上で凝縮し，得られた液体の水はその後冷却しないものとする。凝縮で放出される熱量 $Q_{\text{rel}}$ を表し，$m_v=1{,}0\times10^{-2}\,\text{kg}$ の場合を計算せよ。ガラスの温度では $L_v=2{,}45\times10^6\,\text{J}\!\cdot\!\text{kg}^{-1}$ とする。
\end{enumerate}

\begin{indication}
相転移を伴わない加熱には $Q=mc\Delta T$，一定温度での相転移には $Q=mL$ を用いる。高さ $h$ にある質量 $m$ の重力位置エネルギーは $E_p=mgh$ であり，$g=9{,}81\,\text{m}\!\cdot\!\text{s}^{-2}$ とする。
\end{indication}

\begin{solution}
\textbf{1.} 氷を融点まで加熱するには
\[
Q_1=mc_{\text{ice}}\Delta T
=1{,}0\times2{,}1\times10^3\times100=2{,}1\times10^5\,\text{J}.
\]
\textbf{2.} 融解中，温度は $0\,{}^\circ\text{C}$ のままである：
\[
Q_2=mL_f=1{,}0\times3{,}34\times10^5=3{,}34\times10^5\,\text{J}.
\]
\textbf{3.} 液体の水を加熱するには
\[
Q_3=mc_{\text{water}}\Delta T
=1{,}0\times4{,}18\times10^3\times100=4{,}18\times10^5\,\text{J}.
\]
\textbf{4.} 完全な気化にはさらに
\[
Q_4=mL_v=1{,}0\times2{,}26\times10^6=2{,}26\times10^6\,\text{J}
\]
が必要である。この段階だけで数メガジュールのエネルギーを要する。

\textbf{5.} 気化が圧倒的に最も多くのエネルギーを要する：
\[
Q_4=2{,}26\times10^6>Q_3=4{,}18\times10^5>Q_2=3{,}34\times10^5>Q_1=2{,}1\times10^5\,\text{J}.
\]
気化には，液体の水を $0$ から $100\,{}^\circ\text{C}$ まで加熱する場合の約 $(2{,}26\times10^6)/(4{,}18\times10^5)\approx5{,}4$ 倍のエネルギーが必要である。

$h=10\,\text{m}$ から落下する質量 $m$ について $E_p=mgh$ である。$mgh=Q_3$ とおけば
\[
m=\frac{Q_3}{gh}=\frac{4{,}18\times10^5}{9{,}81\times10}
\approx4{,}26\times10^3\,\text{kg}.
\]
したがって，1キログラムの水を $0$ から $100\,{}^\circ\text{C}$ まで加熱するのと同じエネルギーを放出するには，約 $4{,}3$ トンの質量を10メートル落下させる必要がある。これは非常に大きく，給湯が家庭のエネルギー消費の大きな割合を占める理由を示している。また水の比熱容量は一般的な金属に比べて大きく，例えば銅の約10倍である（後の演習を参照）。

\textbf{6.} 凝縮は気化の逆過程であり，蒸気はガラスに潜熱を放出する：
\[
Q_{\text{rel}}=m_vL_v.
\]
$m_v=1{,}0\times10^{-2}\,\text{kg}$ のとき，
\[
Q_{\text{rel}}=1{,}0\times10^{-2}\times2{,}45\times10^6
=2{,}45\times10^4\,\text{J}.
\]
少量の蒸気の凝縮でも無視できない熱量を放出し，それをフロントガラスと周囲が受け取る。
\end{solution}
\end{exo}

\begin{exo}[大きな樹木の蒸発散：天然のエアコン？]{odg-evapotranspiration-arbre}
\seoready{true}
\keywords{オーダー, 蒸発散, 樹木, 潜熱, 冷却能力, エアコン, COP}

暑く乾燥した日に，十分な水を得られる大きな樹木は，蒸発散によって約 $500\,\text{L}=5{,}0\times10^{-1}\,\text{m}^3$ の水を放出することがある。蒸散は主に光がある間に起こるので，12時間継続すると仮定する。樹冠が地面を覆う面積を $A_{\text{tree}}=50\,\text{m}^2$ とする。水の密度と大気圧下での気化潜熱は
\[
\rho_{\text{water}}=1{,}0\times10^3\,\text{kg}\!\cdot\!\text{m}^{-3},
\qquad L_v=2{,}45\times10^6\,\text{J}\!\cdot\!\text{kg}^{-1}.
\]
\begin{enumerate}
  \item 一日に蒸発散する水の質量と，その水を気化させるのに必要なエネルギーを求めよ。
  \item この過程が冷却効果を生む理由を説明せよ。
  \item 活発な蒸散が続く12時間について，樹木の1平方メートル当たりの平均冷却能力を求めよ。
  \item 比較のため，入力電力 $P_{\mathrm{el}}=2{,}0\times10^3\,\text{W}$，冷房成績係数 $\mathrm{COP}=3{,}0$ で，面積 $A_{\text{room}}=20\,\text{m}^2$ の部屋を冷やすエアコンを考える。定義は $\mathrm{COP}=P_{\text{cool}}/P_{\mathrm{el}}$ である。冷却能力と単位面積当たりの冷却能力を求め，樹木と比較せよ。
  \item 単位面積当たりの能力が同じオーダーでも，樹木とエアコンが全く同じ体感冷却を生むとはいえないのはなぜか。
  \item 多数の観葉植物で家庭を冷やせるだろうか。（一般知識：答えは生物学的過程に依存する。）
\end{enumerate}
\begin{indication}
$m=\rho V$，$Q_{\mathrm{evap}}=mL_v$，$P=Q/\tau$ を用いる。エアコンには $P_{\text{cool}}=\mathrm{COP}\,P_{\mathrm{el}}$ を用いる。
\end{indication}
\begin{solution}
\textbf{1.} 水の質量は
\[
m=\rho_{\text{water}}V=1{,}0\times10^3\times5{,}0\times10^{-1}
=5{,}0\times10^2\,\text{kg}.
\]
気化に必要なエネルギーは
\[
Q_{\mathrm{evap}}=mL_v=5{,}0\times10^2\times2{,}45\times10^6
=1{,}225\times10^9\,\text{J}.
\]
したがってオーダーは1ギガジュールである。

\textbf{2.} 気化は吸熱相転移であり，葉，土壌，周囲の空気からエネルギーを奪う。この熱の除去により葉とその近傍の温度が下がる。

\textbf{3.} 12時間は $\tau=12\times3600=4{,}32\times10^4\,\text{s}$ である。よって
\[
P_{\text{tree}}=\frac{Q_{\mathrm{evap}}}{\tau}
=\frac{1{,}225\times10^9}{4{,}32\times10^4}\approx2{,}84\times10^4\,\text{W},
\]
単位面積当たりでは
\[
\frac{P_{\text{tree}}}{A_{\text{tree}}}
=\frac{2{,}84\times10^4}{50}\approx5{,}7\times10^2\,\text{W}\!\cdot\!\text{m}^{-2}.
\]
\textbf{4.} エアコンの有効冷却能力は
\[
P_{\text{cool}}=\mathrm{COP}\,P_{\mathrm{el}}
=3{,}0\times2{,}0\times10^3=6{,}0\times10^3\,\text{W}.
\]
単位床面積当たりでは
\[
\frac{P_{\text{cool}}}{A_{\text{room}}}
=\frac{6{,}0\times10^3}{20}=3{,}0\times10^2\,\text{W}\!\cdot\!\text{m}^{-2}.
\]
このモデルでは，樹木の単位面積当たり蒸発散能力は一般的なエアコンのほぼ2倍である。

\textbf{5.} 比較しているのはエネルギー流だけである。エアコンは閉じた制御空間を冷やすが，樹木の水蒸気は大気中に拡散し，風は暖気を運ぶ。体感は換気，湿度，気温，樹木からの距離に依存する。蒸散も日射，風，湿度，樹種，利用可能な水によって変わり，干ばつ時には大幅に低下しうる。一方，木陰は地面や人体が受ける日射を直接減らす。したがって計算はオーダーを比較するもので，熱的快適性の厳密な等価性を示さない。

\textbf{6.} 実際には観葉植物の冷却効果はごく小さい。多くの植物では光によって気孔が開き，葉の小さな孔から水蒸気が出る。室内照明は一般に弱いため，気孔の開きと蒸散は減少する。換気の悪い部屋では，湿度の上昇も徐々に蒸発を遅くする。

多数の植物は局所的にわずかな蒸発冷却を生じうるが，エアコンの代わりにはならない。住居から継続的にエネルギーを除くには，湿った空気を屋外へ排出する必要がある。強い人工照明は蒸散を促進しうるが，その電気エネルギーはほぼ全て室内の熱になる。サボテンを含む乾燥環境に適応した CAM 植物は例外で，主に夜間に気孔を開く。
\end{solution}
\end{exo}

\begin{exo}[二つの水塊の混合]{calorimetrie-melange-eau}
\seoready{true}
\keywords{熱量測定, 混合, 熱容量, 比熱容量, ジョゼフ・ブラック, 平衡温度}

$T_1=80\,^\circ\text{C}$ の水 $m_1=0{,}200\,\text{kg}$ を，$T_2=15\,^\circ\text{C}$ の水 $m_2=0{,}300\,\text{kg}$ が入った容器に注ぐ。容器は理想的な熱量計で，外部への熱損失と容器自身の熱容量を無視する。$Q=mc\Delta T$ であり，液体水の比熱容量は $c=4{,}18\times10^3\,\text{J}\!\cdot\!\text{kg}^{-1}\!\cdot\!\text{K}^{-1}$ である。
\begin{enumerate}
  \item 温水が放出した熱を冷水が全て受け取ることを説明し，$m_1$，$m_2$，$c$，$T_1$，$T_2$，平衡温度 $T_{eq}$ を含む保存式を書け。
  \item $T_{eq}$ の式を導き，数値を求めよ。
  \item 混合中に温水が放出した熱量 $Q$ を求めよ。
  \item 同じ物質同士の混合から $c$ を求められるか。理由を述べよ。
\end{enumerate}
\begin{indication}
熱量計は断熱かつ剛体なので，全内部エネルギー $U_1+U_2$ は保存される。一方が放出した熱を他方が逆符号で受け取る。
\end{indication}
\begin{solution}
\textbf{1.}
\[
m_1c(T_{eq}-T_1)+m_2c(T_{eq}-T_2)=0.
\]
\textbf{2.} 共通因子 $c$ が消去され，
\[
T_{eq}=\frac{m_1T_1+m_2T_2}{m_1+m_2}
=\frac{0{,}200\times80+0{,}300\times15}{0{,}200+0{,}300}=41\,^\circ\text{C}.
\]
\textbf{3.}
\[
Q=m_1c(T_1-T_{eq})=0{,}200\times4{,}18\times10^3\times(80-41)
\approx3{,}26\times10^4\,\text{J}.
\]
冷水が受け取る熱量も同じである：
\[
m_2c(T_{eq}-T_2)=0{,}300\times4{,}18\times10^3\times26
\approx3{,}26\times10^4\,\text{J}.
\]
\textbf{4.} 求められない。同じ物質では $c$ が収支式の両項に掛かって消去される。平衡温度は $c$ に依存せず，質量で重み付けした初期温度の平均である。混合法で比熱容量を測るには，次の演習のように既知の熱容量をもつ物質が必要である。
\end{solution}
\end{exo}

\begin{exo}[混合法による金属の比熱容量の決定]{calorimetrie-methode-melanges}
\seoready{true}
\keywords{熱量測定, 混合法, 比熱容量, 熱量計, 水当量}

18世紀のジョゼフ・ブラックと同様に，\emph{混合法}で未知金属の比熱容量 $c_m$ を測定する。質量 $m=0{,}150\,\text{kg}$ の試料を $T_m=95\,^\circ\text{C}$ まで加熱し，$m_e=0{,}250\,\text{kg}$ の水を含む熱量計に素早く入れる。水について $c_e=4{,}18\times10^3\,\text{J}\!\cdot\!\text{kg}^{-1}\!\cdot\!\text{K}^{-1}$，$T_e=18\,^\circ\text{C}$ であり，平衡温度は $T_{eq}=22\,^\circ\text{C}$ である。まず容器，撹拌棒，温度計の熱容量を無視する。
\begin{enumerate}
  \item 金属と水からなる孤立系のエネルギー保存式を，$c_m$ のみを未知数として書け。
  \item $c_m$ の式と数値を求めよ。どの一般的な金属に近いか。基準値の単位は $\text{J}\!\cdot\!\text{kg}^{-1}\!\cdot\!\text{K}^{-1}$ で，アルミニウム $9{,}0\times10^2$，鉄 $4{,}5\times10^2$，銅 $3{,}9\times10^2$，鉛 $1{,}3\times10^2$ である。
  \item 容器も熱を吸収する。この効果を\emph{水当量} $\mu=1{,}5\times10^{-2}\,\text{kg}$ で表し，熱量計を追加の水質量 $\mu$ とみなす。$c_m$ を再計算し，補正の向きを説明せよ。
  \item 試料を素早く入れ，熱量計を周囲の空気から十分断熱する必要があるのはなぜか。
\end{enumerate}
\begin{indication}
金属は $Q_m=mc_m(T_m-T_{eq})$ を放出し，水と，第3問では熱量計が全て受け取る：$Q_m=(m_e+\mu)c_e(T_{eq}-T_e)$。
\end{indication}
\begin{solution}
\textbf{1.}
\[
mc_m(T_m-T_{eq})=m_ec_e(T_{eq}-T_e).
\]
\textbf{2.}
\[
c_m=\frac{m_ec_e(T_{eq}-T_e)}{m(T_m-T_{eq})}
=\frac{0{,}250\times4{,}18\times10^3\times(22-18)}{0{,}150\times(95-22)}
\approx3{,}82\times10^2\,\text{J}\!\cdot\!\text{kg}^{-1}\!\cdot\!\text{K}^{-1}.
\]
これは銅の値に非常に近い。

\textbf{3.}
\[
c_m=\frac{(m_e+\mu)c_e(T_{eq}-T_e)}{m(T_m-T_{eq})}
=\frac{(0{,}250+0{,}015)\times4{,}18\times10^3\times4}{0{,}150\times73}
\approx4{,}05\times10^2\,\text{J}\!\cdot\!\text{kg}^{-1}\!\cdot\!\text{K}^{-1}.
\]
補正により $c_m$ は大きくなる。容器を無視すると金属が実際に放出した熱量を過小評価し，その比熱容量も過小評価するからである。

\textbf{4.} この方法は測定中ずっと系が孤立していることを前提とする。素早い投入は試料が水に達する前の空気との熱交換を抑え，良い断熱は平衡に達する間の損失を抑える。未計上の伝熱は収支と $c_m$ を誤らせる。これはブラック，さらに氷熱量計を用いたラヴォアジエとラプラスが重視した実験上の注意である。
\end{solution}
\end{exo}

\begin{exo}[食品カロリーと代謝出力]{calorimetrie-calories-alimentaires}
\seoready{true}
\keywords{カロリー, キロカロリー, 食品カロリー, 仕事当量, 代謝出力, 単位}

栄養表示では，成人は平均して一日約 $2000\,\text{kcal}$ を摂取するとされる。キロカロリーは栄養学でなお使われる非SI単位であり，$1\,\text{kcal}=4{,}18\times10^3\,\text{J}$ である。
\begin{enumerate}
  \item 一日の摂取量をジュールに換算せよ。
  \item エネルギーを24時間にわたり一様に使うとして，平均代謝出力をワットで求めよ。
  \item エネルギーバーに「$210\,\text{kcal}$」と表示されている。このエネルギーが全て熱に変換されるなら，初温 $20\,^\circ\text{C}$ の水を何キログラム $100\,^\circ\text{C}$ の沸点まで加熱できるか。$c_{\text{water}}=4{,}18\times10^3\,\text{J}\!\cdot\!\text{kg}^{-1}\!\cdot\!\text{K}^{-1}$ とする。
  \item 人体は実際にこのエネルギーをどのような形で利用するか。定性的に答えよ。
\end{enumerate}
\begin{indication}
第2問では $\mathcal P=E/\Delta t$ を用いる。第3問ではまずジュールに換算し，$Q=mc\Delta T$ から $m$ を求める。
\end{indication}
\begin{solution}
\textbf{1.}
\[
E=2000\times4{,}18\times10^3=8{,}36\times10^6\,\text{J}.
\]
\textbf{2.}
\[
\mathcal P=\frac{8{,}36\times10^6}{8{,}64\times10^4}\approx97\,\text{W}.
\]
一日全体では，人間は白熱電球と同程度の平均出力を使う。この意外なオーダーは，食品カロリーを身近な物理単位に換算する意義を示す。

\textbf{3.}
\[
Q=210\times4{,}18\times10^3\approx8{,}78\times10^5\,\text{J}.
\]
$\Delta T=80\,\text{K}$ より，
\[
m=\frac{Q}{c\Delta T}
=\frac{8{,}78\times10^5}{4{,}18\times10^3\times80}
\approx2{,}63\,\text{kg}.
\]
したがって，エネルギーバー1本は，オーダーとして約 $2{,}5\,\text{kg}$ の水道水を沸騰させるだけのエネルギーをもつ。

\textbf{4.} 人体はエネルギーを消滅させるのではなく，栄養素の化学エネルギーを変換する。一部は筋肉の仕事，臓器の機能，細胞の電気的勾配の維持，分子輸送，化学合成に使われ，一部はグリコーゲンや脂肪として蓄えられる。最終的に使用されたエネルギーの大半は熱として散逸し，体温維持に寄与した後，環境へ移る。少量は排泄物に残る化学エネルギーとして体外へ出る。
\end{solution}
\end{exo}
